\emph{Numbers in this section are from the \texttt{quick} experiment profile (Section~\ref{sec:appendix-profiles}); the \texttt{full} profile matches the factorial sizes in our project plan and is intended for an offline run.}

\subsection{RQ1: solver behavior}
Figure~\ref{fig:e01} shows the fitted convergence order for each solver against the gold standard.
Observed orders match theory closely: Forward Euler $\hat p=1.03$ (theory 1), Heun $\hat p=1.94$
(theory 2), classical RK4 $\hat p=3.94$ (theory 4), backward Euler $\hat p=0.97$ (theory 1), and
trapezoidal $\hat p=2.00$ (theory 2).

\begin{figure}[t]
  \centering
  \includegraphics[width=\linewidth]{../figures/e01_convergence.pdf}
  \caption{Observed convergence order matches theory (E01, $R_0=3$ regime).}
  \label{fig:e01}
\end{figure}

Figure~\ref{fig:e02} demonstrates the deliberate finding of Section~\ref{sec:methods}: mass error
stays at roundoff magnitude regardless of solver or step size, while the phase-invariant drift
correctly separates solver quality.

\begin{figure}[t]
  \centering
  \includegraphics[width=\linewidth]{../figures/e02_invariants.pdf}
  \caption{Mass conservation is uninformative; phase-invariant drift is not (E02).}
  \label{fig:e02}
\end{figure}

\subsection{RQ2: solver-in-the-loop parameter recovery}
Figure~\ref{fig:e07} is the mechanism behind our central claim: with noise-free synthetic data
(so any deviation from truth is attributable purely to the fitting solver), bias in the recovered
$\hat\beta$ shrinks with $h$ at each solver's own convergence order. At $(\beta,\gamma)=(0.3,0.1)$,
reducing $h$ from $0.25$ to $0.05$ (a factor of 5) shrinks Euler's bias by a factor of $\approx 5.0$
($4.48\times10^{-3} \to 8.93\times10^{-4}$), Heun's by $\approx 24.2$
($5.88\times10^{-5} \to 2.43\times10^{-6}$), and RK4's by $\approx 606$
($7.17\times10^{-9} \to 1.18\times10^{-11}$) -- close to the theoretical $5^1$, $5^2 = 25$, and
$5^4 = 625$.

\begin{figure}[t]
  \centering
  \includegraphics[width=\linewidth]{../figures/e07_recovery_bias.pdf}
  \caption{Solver-induced bias in $\hat\beta$ scales with the solver's convergence order (E07).}
  \label{fig:e07}
\end{figure}

\subsection{RQ3: noise, sampling, and uncertainty}
Figure~\ref{fig:e08} shows RMSE$(\hat\beta)$ growing with noise level $\sigma$, with bootstrap
confidence bands on the RMSE curve itself; at our reference solver/step size, RMSE grows from
$\approx 1.2\times10^{-11}$ (noise-free) to $5.6\times10^{-4}$ at $\sigma=0.02$ and
$2.9\times10^{-3}$ at $\sigma=0.10$.

\begin{figure}[t]
  \centering
  \includegraphics[width=\linewidth]{../figures/e08_noise_degradation.pdf}
  \caption{RMSE$(\hat\beta)$ vs.\ noise level, with bootstrap CI (E08).}
  \label{fig:e08}
\end{figure}

Figure~\ref{fig:e09} shows RMSE$(\hat\beta)$ as the observation window shrinks from 6 to 1.5 times
the peak time (E09). With dense sampling, truncation costs little: RMSE rises by 34\% at
$dt = 0.5$ and 22\% at $dt = 1$. At $dt = 2$ it doubles ($1.5\times10^{-3} \to 3.0\times10^{-3}$),
so coarse sampling and a short window compound. Every window in this grid still contains the peak;
cutting the observations off \emph{before} the peak is a far stronger effect
(Figure~\ref{fig:e11}). Separately (E10), treating $I_0$ as an unknown third parameter rather than
a known constant raises the condition number of the Jacobian from $\kappa_2 \approx 2.2$--$2.6$ to
$\approx 211$ at $I_0 = 1$ and $\approx 1828$ at $I_0 = 20$. It inflates the variance of
$\hat\beta$ by a factor of 19 at $I_0 = 1$ and 13 at $I_0 = 5$, but only 3.5 at $I_0 = 20$
(10 replicates per cell), so the growth in condition number does not carry over to the variance.

\begin{figure}[t]
  \centering
  \includegraphics[width=\linewidth]{../figures/e09_sampling_window.pdf}
  \caption{RMSE$(\hat\beta)$ vs.\ observation window (as a fraction of 6 times the peak time) and
    sampling interval (E09).}
  \label{fig:e09}
\end{figure}

\subsection{Uncertainty quantification}
Figure~\ref{fig:e11} overlays all four UQ routes and compares profile likelihoods on the full vs.\
a window truncated at 0.3 times the peak time. The truncated profile is nearly flat: it varies by
10\% across the $\beta$ grid, against a 65-fold rise for the full window. This is the signature of
non-identifiability before the peak.

\begin{figure}[t]
  \centering
  \includegraphics[width=\linewidth]{../figures/e11_uncertainty.pdf}
  \caption{Four UQ routes agree closely (left); a truncated window flattens the profile likelihood, the visual signature of non-identifiability (right) (E11).}
  \label{fig:e11}
\end{figure}

\subsection{The crossover \texorpdfstring{$\sigma^*$}{σ*}}
Figure~\ref{fig:e13} overlays the deterministic solver bias in $\hat\beta$ against the
noise-induced standard deviation for Forward Euler at $h = 0.1$; they cross at $\sigma^* = 0.058$
(Section~\ref{sec:methods}). At $h = 0.25$ the bias is 2.5 times larger ($4.48\times10^{-3}$, the
same value E07 finds) and $\sigma^* = 0.156$, so $\sigma^*$ grows roughly in proportion to $h$: a
first-order method's bias scales as $h$, while the noise-induced spread grows with $\sigma$. At
$h = 0.5$ the spread is still below the bias at the largest noise level tested ($\sigma = 0.2$), so
there $\sigma^* > 0.2$. For $\hat\gamma$ the crossover comes earlier: $\sigma^* = 0.029$, $0.054$,
and $0.128$ at $h = 0.1$, $0.25$, and $0.5$.

\begin{figure}[t]
  \centering
  \includegraphics[width=\linewidth]{../figures/e13_crossover.pdf}
  \caption{Crossover for Forward Euler at $h = 0.1$: below $\sigma^* = 0.058$ solver bias dominates
    the error in $\hat\beta$; above it, noise dominates (E13).}
  \label{fig:e13}
\end{figure}

\subsection{Extensions: model variants and real data}
The SEIR, SIRS, and SIR-with-vaccination variants (E12) preserve the RQ1 solver ordering
(Euler~$<$~Heun~$<$~RK4 in accuracy at fixed $h$), indicating the finding is not an artifact of the
plain SIR system's particular structure. Applying the same pipeline to the 1666 Eyam plague outbreak
(E14) recovers an $\hat R_0$ broadly consistent with the historical literature; we flag this result
as a demonstration rather than a validated finding, since the digitized case counts we used have not
been independently verified against the primary source (see \texttt{data/raw/provenance.md}).
